\SourceLeaf{336}%
\textsc{Bessel} für \(a\) in Toisen angegeben hat (Art. 5), dem Logarithmus, falls \(r\) ursprünglich in Toisen ausgedrückt war, den folgenden Betrag hinzufügen
\[
0{,}2897827662
\]
oder, falls \(r\) ursprünglich in französischen gesetzlichen Metern angegeben war, vom Logarithmus den folgenden Betrag abziehen
\[
0{,}0000371638
\]

Die Glieder mit den Faktoren (3), (4) usw. lassen sich als Korrekturen auffassen, die aus den angenäherten Logarithmen \(\log\tau,\log\beta,\log\lambda\) die berichtigten Logarithmen \(\log t,\log b,\log l\) machen. Diese Korrekturen sind in allen Fällen, für die unsere Methode bestimmt ist, sehr kleine Dezimalbrüche. Da man jene Logarithmen in der Regel siebenstellig berechnet, ist es bequem, auch die Korrekturen unmittelbar in Einheiten der siebten Dezimalstelle zu erhalten. Dazu setzt man für die Koeffizienten (3), (4) usw. das Zehnmillionenfache der im vorigen Artikel angegebenen Werte ein oder erhöht ihre Logarithmen um sieben Einheiten. So sind sie in unserer Tafel angesetzt. Dementsprechend setzt man auch
\[
\begin{aligned}
\log(7)&=4{,}62872\ (-10)\\
\log3(7)&=5{,}10584\ (-10)
\end{aligned}
\]
Auch in dieser Form sind (3), (4), (5), (6) ebenso wie (1) und (2) noch echte Brüche, ihre Logarithmen also an sich negativ. In der Tafel stehen sie in der üblichen Form, indem für jeden Logarithmus zehn Einheiten geborgt sind.

\Article{35}
Alles, was oben (Art. 20) über die Lösung derselben Aufgabe auf der Kugeloberfläche gesagt wurde, gilt fast unverändert und mit nur geringfügigen Anpassungen auch für die Anwendung unserer Formeln auf die zu Beginn dieser Abhandlung gestellte Aufgabe. Die tatsächlich gegebenen Größen, nämlich die Breite und das Azimut am ersten Ort, bezeichne man mit \(B^0\) und \(T^0\). Zunächst geht man von einem Näherungswert für \(T\) aus (wenn sonst keine Angaben vorliegen, kann man dafür %
